% explainer-source-hash: b61f9ba78915fdf5b2011a630b8496641f14462f328b2f9ba5f40a7d1eb57aaa
% explainer-source-commit: 4302e8d356b88b447a7e8b3cbbfab75e1d7e0dc8
% explainer-generated: 2026-07-16
% Regenerate with /explain-all (see WIRING.md). Do not edit this stamp by hand:
% it is how the toolchain knows whether this document still matches the code.
\documentclass[11pt,a4paper]{article}
\input{preamble}

\title{\vspace{-1.5cm}\textbf{raft-kv}\\[0.3cm]
\large A Distributed Key-Value Store with Raft Consensus,\\
Deterministic Simulation, and a Linearizability Checker\\[0.4cm]
\normalsize Brief: the argument, the architecture, and the evidence}
\author{}
\date{}

\begin{document}
\maketitle
\thispagestyle{empty}

% =====================================================================
\section{The one-paragraph version}

\begin{keyidea}{What this is}
A key-value database that keeps identical copies of its data on several
machines, so that any one of them can die without losing data or serving a
wrong answer. Keeping machines in exact agreement while they crash and their
network breaks is the problem called \textbf{consensus}, and this project
solves it with the \textbf{Raft} algorithm, written from scratch in Python with
no consensus libraries and zero runtime dependencies. But the consensus code is
the smaller half of the project. The larger half is the evidence: a
\textbf{deterministic simulator} that replays thousands of randomized disasters
against the real consensus code as a pure function of a single integer, and a
from-scratch \textbf{linearizability checker} that verifies the answers clients
actually received could have come from a single, honest, non-distributed
database.
\end{keyidea}

\subsection{The thesis}

The project exists to make an argument, and the argument is the interesting part:

\begin{quote}
\itshape Consensus code is easy to write and nearly impossible to trust from
example-based tests, because the bugs live in message interleavings you will
never hit on a laptop loopback. So the system should be testable as a
deterministic function of a seed.
\end{quote}

This is the position taken by the FoundationDB database, and it drives every
structural decision here. If every source of nondeterminism (time, message
delay, loss, duplication, partition timing, crash timing, client behaviour,
election jitter) is drawn from one seeded generator, then a whole 12-second
cluster lifetime under sustained chaos becomes a \textbf{pure function of one
integer}. That buys three things at once:

\begin{itemize}
\item thousands of universes run in seconds, because no real time passes;
\item a failure report is one integer, not a flaky-test shrug;
\item the failing schedule replays exactly, forever, under a debugger.
\end{itemize}

The whole architecture is bent around one precondition for this to be worth
anything: \textbf{the simulated system must be the real system}. A simulator
that tests a copy of the consensus logic tests nothing.

% =====================================================================
\section{Three ideas you need, from first principles}

\begin{jargon}{Consensus}
Getting a group of unreliable machines to agree on one \textbf{ordered list of
decisions}, such that once a decision is made it is permanent, and no two
machines ever disagree about what decision number 7 was. Crucially, the group
does not agree on ``the value of key \code{city}''. It agrees on an ordered log
of commands. If every machine holds the same log and runs the same
deterministic program over it, every machine ends up with identical data,
necessarily. Agreement on the data is derived free from agreement on the order.
This is called a \textbf{replicated state machine}.
\end{jargon}

\begin{jargon}{Quorum, and why it makes split brain impossible}
Every important decision requires a strict majority of the cluster. The reason
this works is a fact about sets, not computers: \textbf{any two majorities of
the same set must share a member}, because 2 + 2 > 3. So a network partition
can leave at most \emph{one} side with a majority; the other side is
mathematically unable to act. Without this, two halves of a split cluster would
each keep accepting writes and diverge permanently, which is called
\textbf{split brain}. In this code the rule is one line:
\code{len(self.config) // 2 + 1}.
\end{jargon}

\begin{jargon}{Linearizability}
The correctness target, and the strongest single-object guarantee. A system is
linearizable if every operation appears to take effect
\textbf{instantaneously at some single moment between the client sending the
request and receiving the response}, consistent with real time. The practical
meaning: if your write finished at 10:00:01, every read starting at 10:00:02,
on any node, must see it. The distributed system becomes indistinguishable from
a single-threaded program holding one copy of the data. It is a property of the
\emph{observable history}, not of the implementation, which is exactly why it
can be checked from the outside without trusting the code.
\end{jargon}

% =====================================================================
\section{The architecture, and the decision it turns on}

\begin{figure}[H]
\centering
\resizebox{0.99\textwidth}{!}{
\begin{tikzpicture}[node distance=9mm]
\node[draw=inkblue, very thick, fill=softgrey, rounded corners=4pt,
      minimum width=50mm, minimum height=26mm] (core) {};
\node[font=\small\bfseries, text=inkblue, anchor=north] at ($(core.north)+(0,-2mm)$)
      {raft/ : the sans-IO core};
\node[box, minimum width=40mm] (rn) at ($(core.center)+(0,1mm)$) {RaftNode};
\node[lbl, anchor=south, align=center] at ($(core.south)+(0,2mm)$)
      {\itshape no asyncio, no time,\\\itshape no socket, no global random};

\node[dec, left=16mm of core, minimum width=17mm] (env) {Env};
\node[dec, right=16mm of core, minimum width=17mm] (sto) {Storage};
\draw[flow] (rn) -- node[lbl, above]{now, call\_later,\\send, rng} (env);
\draw[flow] (rn) -- node[lbl, above]{load, append,\\snapshot} (sto);

\node[extbox, above=13mm of env, minimum width=40mm] (aenv)
  {AsyncioEnv + TcpTransport\\\scriptsize{}real clock, real sockets};
\node[extbox, below=13mm of env, minimum width=40mm] (senv)
  {SimEnv + SimNetwork + SimClock\\\scriptsize{}virtual time, seeded chaos};
\node[extbox, above=13mm of sto, minimum width=40mm] (fs)
  {FileStorage\\\scriptsize{}fsync, CRC, atomic rename};
\node[extbox, below=13mm of sto, minimum width=40mm] (ms)
  {MemoryStorage\\\scriptsize{}a perfect disk};

\draw[dflow] (aenv) -- (env); \draw[dflow] (senv) -- (env);
\draw[dflow] (fs) -- (sto);   \draw[dflow] (ms) -- (sto);

\node[lbl, text=okgreen, anchor=west, align=left] at ($(aenv.east)+(4mm,0)$)
  {\textbf{PRODUCTION}};
\node[lbl, text=accent, anchor=west, align=left] at ($(senv.east)+(4mm,0)$)
  {\textbf{SIMULATION}};
\node[store, minimum width=24mm, below=13mm of core] (kv) {KVMachine\\\scriptsize{}get/put/delete/cas\\\scriptsize{}+ session table};
\draw[flow] (rn) -- node[lbl, right]{apply(entry)} (kv);
\end{tikzpicture}}
\caption{The whole system in one picture. \code{RaftNode} can touch the world
only through two narrow protocols. Swap what is plugged into them and the
identical consensus code runs over real TCP or inside the simulator.}
\end{figure}

\begin{keyidea}{Sans-IO is the whole ballgame}
\textbf{Sans-IO} (``without IO'') means the protocol implementation performs no
input or output itself: it never opens a socket, never reads a clock, never
sleeps. \code{RaftNode} reacts to exactly three stimuli (a message arrived, a
timer fired, the local application asked for something) and emits effects
through exactly two protocols, \code{Env} (\code{now}, \code{call\_later},
\code{send}, \code{rng}) and \code{Storage}.

The payoff is the entire project. Because the node takes its randomness from
\code{env.rng} rather than the global \code{random} module, and its time from
\code{env.now()} rather than \code{time.time()}, the simulator can hand it a
seeded generator and a virtual clock and it cannot tell. So ``run 5000 chaos
schedules in a few minutes'' becomes an ordinary function call, and
\textbf{every bug the simulator finds is a production bug}, because it is
literally the production code. There is no test double of the algorithm
anywhere in this repository, and \file{raft/} contains no \code{if testing:}
branch.
\end{keyidea}

The module layout follows directly, with the dependency arrow pointing inward:

\begin{center}
\begin{tabular}{lp{9.3cm}}
\toprule
\textbf{Module} & \textbf{Role} \\
\midrule
\file{raft/} & the consensus core: \code{RaftNode} (576 lines), the log, storage, messages, and the \code{Env} protocol. Imports nothing outward. \\
\file{kv/} & the replicated state machine: get, put, delete, cas, plus the session table \\
\file{transport/} & real asyncio TCP, length-prefixed JSON frames, and \code{AsyncioEnv} \\
\file{sim/} & virtual clock, seeded chaos network, the harness, the batch runner \\
\file{checker/} & history recording and the linearizability search \\
\file{server.py}, \file{client.py}, \file{cli/} & one TCP node, a retrying client, the \code{raftkv} command \\
\bottomrule
\end{tabular}
\end{center}

% =====================================================================
\section{Raft, in one page}

Every node is a follower, a candidate, or a leader. Time is divided into
\textbf{terms}, an ever-increasing logical clock; at most one leader exists per
term, and any message carrying a higher term instantly demotes its receiver to
follower.

\begin{figure}[H]
\centering
\resizebox{0.96\textwidth}{!}{
\begin{tikzpicture}
\node[box, minimum width=28mm, minimum height=12mm] (f) {\textbf{FOLLOWER}};
\node[box, minimum width=28mm, minimum height=12mm, right=48mm of f] (c) {\textbf{CANDIDATE}};
\node[box, minimum width=28mm, minimum height=12mm, right=48mm of c] (l) {\textbf{LEADER}};

\draw[flow] (f) -- node[lbl, above, align=center]
  {\textbf{election timeout}\\no heartbeat for a random\\time in [T, 2T]} (c);
\draw[flow] (c) -- node[lbl, above, align=center]
  {\textbf{votes from a majority}\\$\Rightarrow$ append a no-op entry,\\start heartbeats} (l);
\draw[flow] (c) to[out=205, in=335] node[lbl, below, align=center]
  {discovers a current leader,\\or sees a higher term} (f);
\draw[flow] (l) to[out=250, in=290, looseness=0.6] node[lbl, below, align=center]
  {\textbf{sees a higher term}, or is removed from the config\\
   $\Rightarrow$ step down: cancel heartbeats, fail all pending reads and writes} (f);
\draw[flow] (c) to[out=70, in=110, looseness=5] node[lbl, above, align=center]
  {split vote: nobody won\\$\Rightarrow$ new term, campaign again} (c);
\end{tikzpicture}}
\caption{The role state machine. Note the asymmetry: there is no
follower-to-leader edge and no leader-to-candidate edge. Power is only taken
through an election and only surrendered by stepping down.}
\end{figure}

\textbf{Election.} Followers wait a \emph{randomized} time (uniform in
\code{[T, 2T]}) without hearing a heartbeat, then increment the term, vote for
themselves, and ask for votes. Randomization is what prevents everyone from
campaigning at once and splitting the vote forever. Two rules carry the weight:
the vote is \textbf{persisted with fsync before the reply is sent} (a node that
forgot its vote after a crash could vote twice in one term, producing two
leaders), and voters \textbf{refuse any candidate whose log is less up to date
than their own}. That second rule is what makes committed data survive every
leader change automatically: a winner needs a majority, any majority overlaps
the majority holding a committed entry, and those voters will refuse anyone
missing it.

\textbf{Replication.} The leader appends a command to its own log, fsyncs it,
and broadcasts an \code{AppendEntries} message. That message carries the index
and term of the entry \emph{immediately before} the new ones. If a follower
disagrees there, it rejects and returns a \textbf{conflict hint}, the first
index of the conflicting term, so the leader can back up in one jump rather
than one entry per round trip. Agreement on that one preceding entry proves agreement on the
entire prefix, by induction over the fact that indexes are never reused and one
leader exists per term. Once a majority stores an entry, it is
\textbf{committed}: permanent, present in every future leader's log.

\begin{keyidea}{The subtlest rule in Raft, and the reason for the no-op}
A leader may \textbf{not} commit an entry from an earlier term merely because it
is now replicated on a majority. If it did, a different node whose log is more
up to date on other grounds could still win a later election, not have that
entry, and overwrite it: an entry reported committed would vanish. So the code
checks \code{self.log.term\_at(candidate) == self.current\_term} before
advancing the commit index (the paper's section 5.4.2). To make progress on the
inherited backlog anyway, a fresh leader immediately appends a \code{noop}
entry of its own term; when that commits by the ordinary majority rule,
everything before it commits transitively and safely. The no-op is not a
nicety: without it a leader receiving no writes could never commit its
predecessor's backlog at all.
\end{keyidea}

\textbf{Crash safety} falls out of one discipline: never acknowledge before
\code{fsync}, and make every on-disk structure either atomic-rename recoverable
or scan-and-truncate recoverable. \code{hardstate.json} (term and vote) is
written temp-file, fsync, rename, then \textbf{the directory is fsynced too},
which is the step most implementations forget. \code{log.bin} is append-only
records of \code{[4-byte length][4-byte CRC32][JSON]}; on recovery the file is
scanned and the first short or bad-CRC record marks a torn tail from a crash
mid-write, which is truncated away. Discarding data is provably safe here
precisely because nothing is acknowledged before fsync returns, so a
half-written record was never promised to anyone.

\textbf{Snapshots} replace an applied prefix of the log with an image of the
state machine (KV data \emph{and} the session table), sent to hopelessly lagging
followers via \code{InstallSnapshot}. \textbf{Membership changes} use the
dissertation's single-server approach: add or remove exactly one node per config
entry, so any two configurations share a common majority and disjoint quorums
cannot form, which is why no two-phase joint-consensus protocol is needed.

% =====================================================================
\section{Linearizable reads: ReadIndex}

\begin{keyidea}{A leader cannot know that it is still the leader}
This is the trap that makes reads harder than writes. Suppose n1 is partitioned
away from n2 and n3, which elect a new leader and accept writes. n1 has heard
nothing, so it still believes it leads, and its local data looks perfectly fine
\emph{to it}. It will happily serve a value that is minutes stale. Nothing in
its own state betrays it. Leadership is not a local fact, so reading from local
memory is unsafe.
\end{keyidea}

The fix is \textbf{ReadIndex}: before answering a read, the leader must
\emph{prove} its currency. It records \code{readIndex = commitIndex}, broadcasts
a heartbeat round carrying a sequence number, and waits for acks from a majority
for that round; only then does it read local state. A deposed leader in a
minority partition is not merely unlikely to answer, it is \textbf{structurally
unable} to, because those acks cannot cross the partition. When it eventually
sees the higher term it steps down, and stepping down explicitly fails every
pending read, so the client is told \code{not\_leader} and retries elsewhere.

There is a second precondition, and the README credits the chaos harness with
finding it: a fresh leader must first commit an entry \textbf{of its own term}
before serving reads, because its \code{commit\_index} is a stale local variable
that may be lower than what is actually committed cluster-wide. Reads arriving
too early are parked with no read index and assigned one the instant the
leader's no-op commits.

\begin{keyidea}{The alternative that was refused, and why the reason matters}
\emph{Leader leases} let a leader serve reads for
\code{election\_timeout - clock\_drift} after one heartbeat quorum, saving a
round trip per read. This project refuses them, and documents the refusal:
leases buy latency at the price of a safety argument that depends on
\textbf{bounded clock drift between physical machines}. In a project whose
entire thesis is testable correctness, importing an untestable assumption about
hardware clocks is a bad trade. ReadIndex costs one heartbeat round and assumes
nothing about clocks.
\end{keyidea}

% =====================================================================
\section{Exactly-once writes}

\begin{keyidea}{Why retries corrupt data, and why the fix must live in the replicated state}
A client sends a cas. The leader commits it, then dies before replying. The
client times out and retries against the new leader, which has no idea this is
a retry and applies it again. For a cas the retry now \emph{fails} (the value
already changed), so the client is told its operation failed when in fact it
succeeded. The fix: every write carries \code{(client\_id, seq)}, and the state
machine remembers per client the last seq applied and \textbf{the result it
returned}; a duplicate returns the cached result without re-executing.

The part that is easy to get wrong is where the table lives. It must be
\textbf{inside the replicated state}, going through the log and riding inside
snapshots. A leader-local cache would silently downgrade the system from
exactly-once to at-least-once precisely at failover, which is exactly when it
matters. Exactly-once is a state machine feature, not a networking feature.
\end{keyidea}

% =====================================================================
\section{The evidence machine}

\subsection{One seeded run}

\begin{figure}[H]
\centering
\resizebox{0.99\textwidth}{!}{
\begin{tikzpicture}[node distance=8mm]
\node[box, minimum width=22mm] (seed) {\textbf{seed}};
\node[dec, right=11mm of seed, minimum width=15mm] (root) {root\\Random};
\draw[flow] (seed) -- (root);
\node[extbox, right=12mm of root, minimum width=27mm, yshift=13mm] (r1) {network rng};
\node[extbox, right=12mm of root, minimum width=27mm, yshift=0mm] (r2) {per-node rngs};
\node[extbox, right=12mm of root, minimum width=27mm, yshift=-13mm] (r3) {nemesis + client rngs};
\draw[dflow] (root) -- (r1); \draw[dflow] (root) -- (r2); \draw[dflow] (root) -- (r3);

\node[box, minimum width=28mm, minimum height=26mm, right=12mm of r2] (heap)
  {\textbf{SimClock}\\one heap ordered\\by (time, seq)};
\draw[flow] (r1) -- (heap); \draw[flow] (r2) -- (heap); \draw[flow] (r3) -- (heap);

\node[box, minimum width=26mm, right=12mm of heap, yshift=9mm] (nodes) {3 x RaftNode\\\scriptsize{}the real thing};
\node[box, minimum width=26mm, right=12mm of heap, yshift=-9mm] (mon) {Monitor\\\scriptsize{}4 live invariants};
\draw[flow] (heap) -- (nodes);
\draw[flow] (nodes) -- node[lbl, right]{trace} (mon);
\node[box, minimum width=26mm, right=13mm of nodes] (chk) {linearizability\\checker};
\draw[flow] (nodes) -- node[lbl, above]{history} (chk);
\node[box, minimum width=22mm, right=11mm of mon] (out) {PASS / FAIL\\\scriptsize{}named by seed};
\draw[flow] (mon) -- (out);
\draw[flow] (chk) to[out=0, in=90] (out);
\end{tikzpicture}}
\caption{Everything derives from the seed; everything funnels through one
ordered event heap; two independent oracles judge the result. The tie-break on
a sequence number inside that heap is load-bearing: without it, two events at
the same virtual instant would run in an undefined order and determinism would
die.}
\end{figure}

Each seed runs 12 virtual seconds of a 3-node cluster with 3 concurrent clients
on 3 keys, under 5\% message loss, 5\% duplication, randomized delay, and a
nemesis firing at Poisson intervals: random bisecting partitions, crash with
delayed restart, election-timer clock skew (0.5x to 1.8x), and heals. Because
time is a variable rather than something you wait for, those 12 virtual seconds
cost roughly 0.1 real seconds. There is no \code{sleep} anywhere in the
simulator.

Chaos stops at 70\% of the run and the cluster is then healed and stabilized so
in-flight operations can drain. That is not cheating, it is what makes the check
bite: \textbf{incomplete operations are the easiest thing for a linearizability
checker to accept}, because it can always declare they never happened. A run
that ended mid-chaos would look clean while testing very little.

Four invariants are checked live, each failure naming its seed: election safety
(at most one leader per term, tracked across the whole run), leader completeness
(a new leader's log is checked against everything ever applied, at the moment it
wins), state machine safety (two nodes applying different entries at one index
is an immediate failure), and log matching.

\subsection{The linearizability checker}

Written from scratch: a Wing \& Gong exhaustive search with Lowe's memoization.
Repeatedly pick an operation that is \emph{minimal} (real-time order permits it
to go next), apply it to a model of a single register, and check the model
reproduces what the client actually observed; backtrack when it cannot. Two
decompositions make it tractable: keys are independent registers so each key is
checked separately, and the set of linearized operations is a \textbf{bitmask
integer}, making the memo key a single machine word. Incomplete operations
branch both ways, implemented elegantly as a weaker acceptance test (every
\emph{completed} op must be linearized; leftovers may be declared never to have
happened).

\begin{keyidea}{Three outcomes, never two}
\code{OK}, \code{VIOLATION}, and \code{UNKNOWN}. The third is the honest one.
Linearizability checking is NP-complete in general, so any checker must
eventually refuse; this one refuses explicitly above 63 operations per key or
2,000,000 explored states. The property that matters: it \textbf{never returns a
false PASS or a false FAIL}, only sometimes a shrug. Note also that exhausting
the search space with no explanation found is a \emph{proof} of violation, not a
timeout.
\end{keyidea}

\subsection{Testing the tests}

\begin{keyidea}{A safety net that never fires is indistinguishable from one that cannot fire}
This is the best idea in the suite, and it appears three times independently.
\code{test\_broken\_raft\_is\_caught} monkeypatches the quorum rule from
\code{// 2 + 1} to \code{// 2}, deliberately making split brain possible, and
asserts the harness detects a violation within 40 seeds. The checker is fed
known-bad synthetic histories (stale read, lost update, phantom value) and must
reject them. And \code{test\_same\_seed\_same\_run} runs a seed twice and
compares a full history fingerprint, so determinism is asserted structurally
rather than assumed.
\end{keyidea}

The README tells one story on itself that is worth repeating, because it is the
same lesson at a smaller scale. The first ``checker budget'' test used 40
incomplete concurrent puts and passed, but it passed for the wrong reason: a
history of purely incomplete operations is trivially linearizable by declaring
none of them happened, so the search finished instantly and exercised none of
the mechanism it claimed to test. The fixed test buries a real violation under
20 \emph{completed} concurrent puts, forcing the exhaustive search the budget
exists to bound. A test that passes tells you nothing until you know why.

% =====================================================================
\section{Verified, not quoted}

The claims below were \textbf{executed on this machine} for this document, from
a clean virtual environment (Linux, Python 3.12.3, 8 cores), rather than copied
from the README.

\begin{center}
\begin{tabular}{lll}
\toprule
\textbf{Claim (README)} & \textbf{Observed} & \textbf{Verdict} \\
\midrule
57 tests, all passing & \code{57 passed in 51.41s} & confirmed \\
5000 seeds swept & 5000 seeds & confirmed \\
\code{564170} client ops total & \code{564170} client ops total & confirmed, exactly \\
0 failures & \code{0 failure(s)} & confirmed \\
84.7s (59.0 seeds/s), 10 workers & 257.0s (19.5 seeds/s), 8 workers & differs; see below \\
\bottomrule
\end{tabular}
\end{center}

\begin{lstlisting}[caption={The headline sweep, run for this document}]
$ .venv/bin/raftkv sim --seeds 5000 --workers 8
5000 seeds in 257.0s (19.5 seeds/s), 564170 client ops total, 0 failure(s)
\end{lstlisting}

\begin{keyidea}{The op count reproduced to the digit, and that is the real result}
The README's run used 10 workers; this one used 8, on a different machine, on a
different day. The operation count came back \textbf{identical}. That is not a
coincidence and it is not a weak check: it is the determinism claim confirmed
end to end. Had any part of the system consulted a wall clock, a hash seed, a
set iteration order, or the global \code{random} module, the same 5000 seeds
would have produced a different number. They did not.
\end{keyidea}

\begin{honest}{The timing did not reproduce, and the gap is larger than the worker count explains}
84.7s at 10 workers versus 257.0s at 8 workers is a factor of about 3; the
worker ratio accounts for at most 1.25. This is not a discrepancy in the code:
throughput is deliberately \emph{outside} the seed and depends on the machine
and its load. \textbf{This run shared the machine with other substantial
workloads}, which is the most likely explanation, but that is an
\textbf{inference}: it was not isolated or measured, and the original run was
not repeated under matched conditions. The honest position: the timing figure is
a machine-and-day observation rather than a property of the project, while the
two figures that \emph{are} properties of the project, the op count and the
failure count, both reproduced exactly.
\end{honest}

\begin{honest}{What the headline number does and does not cover}
``5,000 randomly generated disasters and 564,170 reads/writes with 0
contradictions'' is true and was reproduced. Its scope is narrower than it
sounds, and saying so is what makes it defensible:
\begin{itemize}
\item \textbf{5000 seeds is a deep search of one fixed scenario}, not a broad
      search across scenarios: every seed runs the same recipe (3 nodes, 12
      virtual seconds, the same fault probabilities) and varies only timing,
      targets and ordering.
\item \textbf{The ops are get, put and cas.} The simulated client has no
      \code{delete} branch, so \code{delete} is never exercised end to end by
      any of the 5000 seeds.
\item \textbf{The simulated disks are perfect.} The simulator uses
      \code{MemoryStorage}, an ideal disk that never tears a write. The
      excellent torn-tail CRC recovery code in \code{FileStorage} is therefore
      never touched by the sweep; it is covered only by targeted unit tests.
      This is the largest gap between what the simulator tests and what
      production runs, and the README names it.
\item \textbf{Membership changes are never exercised by the sweep} either: the
      nemesis crashes, partitions and skews, but never adds or removes a server.
\item \textbf{``0 contradictions'' means no violation fired}, not that every
      history was proven linearizable. \code{SimResult.ok} accepts
      \code{UNKNOWN}, and the batch runner never counts it, so the sweep cannot
      distinguish 5000 proofs from 4990 proofs and 10 shrugs. The configured
      width almost certainly keeps every history inside the budget, but the
      runner does not \emph{show} that, and it easily could.
\end{itemize}
\end{honest}

% =====================================================================
\section{Honest assessment}

\subsection{What is genuinely strong}

\begin{itemize}
\item \textbf{The sans-IO boundary is real, not aspirational.} \file{raft/}
      imports \code{enum}, \code{typing} and its own siblings, and nothing else.
      The central architectural claim survives a grep.
\item \textbf{The subtle rules are all present and correct}: vote persistence
      before reply, the up-to-date check, the section 5.4.2 own-term commit
      restriction with its no-op, a conflict hint that can only move
      \code{next\_index} downward (so a confused hint costs a round trip, never
      correctness), the directory fsync after rename, config effective on
      append, and a leader that defers its own removal until it commits.
\item \textbf{The tests test the tests}, three independent ways.
\item \textbf{The honesty is load-bearing}: the \code{UNKNOWN} outcome, the
      documented rejection of leader leases with reasons, and the README's own
      limits section are the work of someone optimizing for being right rather
      than for looking right.
\end{itemize}

\subsection{Weaknesses, ranked by what would hurt most under questioning}

\begin{enumerate}
\item \textbf{The simulator's disk is perfect, so the crash-recovery code is the
      least-tested important code in the project.} The project's thesis is
      ``search the schedule space''; storage faults are simply outside the
      searched space. This is the answer to ``what would you do next'' and
      should be volunteered, not conceded.
\item \textbf{\code{UNKNOWN} is invisible} in \code{SimResult.ok} and in the
      batch runner's output. A few lines would make the headline strictly
      stronger.
\item \textbf{Replication is naive and worst-case quadratic.} Every heartbeat to
      a lagging follower re-sends the \emph{entire} remaining suffix, because
      \code{next\_index} only advances on a reply. Correct, since AppendEntries
      is idempotent, but with a 30 ms heartbeat and a follower 10,000 entries
      behind that is the whole suffix serialized to JSON every 30 ms until the
      first reply lands. There is no batch-size cap.
\item \textbf{\code{\_load\_snapshot()} re-reads and re-parses the entire
      persistent state, including the whole log, on every InstallSnapshot
      send}, then ships the snapshot whole with no chunking.
\item \textbf{No pre-vote}, so a healed node forces one pointless re-election
      through term inflation.
\item \textbf{Sessions never expire}, and the table is inside the replicated
      state and every snapshot. Since the CLI builds a fresh client with a fresh
      random id per invocation, every \code{raftkv put} from the command line
      permanently adds a row.
\item \textbf{Scale}: Python, JSON on the wire, one fsync per append batch, one
      TCP connection per client request. The README's first honest limit says
      plainly that this is a correctness project and not a fast store, which is
      the right call.
\end{enumerate}

\subsection{Smaller things a careful reader will find}

\begin{itemize}
\item The checker validates \code{get} and \code{cas} results against the model
      but treats \code{put} and \code{delete} as pure state transitions, so
      \code{delete}'s \code{existed} flag is never checked.
\item The log-matching invariant breaks at the first index where terms agree,
      inferring the prefix from the Log Matching Property, which is the property
      under test. It is partially circular; other layers still catch what it
      misses.
\item \code{next\_index} initialization disagrees with itself across three
      functions (off by one in \code{\_change\_config}); self-correcting, so the
      tests never see it.
\item \code{\_config\_from}'s index filter is effectively vacuous at both call
      sites.
\end{itemize}

% =====================================================================
\section{What to take away}

\begin{keyidea}{Five sentences}
\textbf{One.} Consensus reduces to keeping an ordered log identical across
machines; Raft does it by electing one leader per term who alone appends, and by
requiring a majority for every decision, so that the overlap between any two
majorities makes contradiction impossible.

\textbf{Two.} Everything rests on sans-IO: \code{RaftNode} touches the world
only through \code{Env} and \code{Storage}, so identical code runs over TCP and
under the simulator, and every simulator bug is by construction a production
bug.

\textbf{Three.} Because all nondeterminism funnels through one seeded generator
and one ordered event heap, a chaotic cluster lifetime is a pure function of one
integer: thousands of disasters searchable in minutes, every failure
reproducible forever.

\textbf{Four.} Correctness is judged from the outside by a from-scratch
linearizability checker that proves, refutes, or honestly says \code{UNKNOWN},
and never guesses.

\textbf{Five.} The safety nets are themselves tested, because a net that never
fires is indistinguishable from one that cannot.
\end{keyidea}

\begin{center}
\begin{tabular}{p{5.3cm}p{8.9cm}}
\toprule
\textbf{Likely question} & \textbf{Honest answer} \\
\midrule
Why not use an existing Raft library? & The point was to learn where the subtlety lives. Section 5.4.2 and the ReadIndex term-commit precondition are two places a plausible shortcut is a real linearizability bug, and the chaos harness found the second one during development. \\
Is 5000 seeds meaningful? & It is a deep search of one fixed scenario, not a broad one. It reproduced to the digit at a different worker count, which is the determinism claim confirmed. \\
What is not tested? & Storage faults (the simulated disk is perfect), \code{delete} end to end, membership changes under chaos, and anything beyond single-digit node counts on localhost. \\
What would you do first? & Inject torn writes and fsync reordering into simulated storage, so the sweep covers the recovery paths that are currently only unit-tested. Then report \code{UNKNOWN} counts. \\
Why ReadIndex, not leases? & Leases trade a round trip for an assumption about bounded clock drift between machines. In a project about testable correctness, an untestable hardware assumption is a bad trade. \\
\bottomrule
\end{tabular}
\end{center}

\vfill
\begin{center}
\small\itshape
Every figure in the verification section was produced by running the code on
this machine, not by quoting the README. Anything inferred is labelled inferred.
Where code and documentation disagreed, the code was taken as the truth.
For the exhaustive treatment, including the full state machine, the election and
replication paths, the partition analysis, and a module-by-module walkthrough,
see the companion Deep Dive.
\end{center}

\end{document}
